\documentclass[11pt,a4paper]{article}
\usepackage{xeCJK}
\usepackage{fontspec}
\usepackage{xcolor}
\usepackage{geometry}
\usepackage[protrusion=true]{microtype}
\geometry{margin=1.8cm,top=1.9cm}
\linespread{1.25}

\setmainfont{Baskerville}
\setCJKmainfont{Songti SC}

\definecolor{tipblue}{RGB}{0,95,160}
\definecolor{rulegray}{RGB}{190,190,190}

\setlength{\parindent}{0pt}
\setlength{\parskip}{0.35em}
\newcommand{\num}[1]{\makebox[1.9em][r]{#1.}\hspace{0.3em}}
\newcommand{\wline}{\noindent\rule{\linewidth}{0.4pt}}
% 中文题干：宋体粗体；提示词：蓝色小字；英文句子：Baskerville 正体
\newcommand{\prompt}[1]{\emph{\textcolor{tipblue}{\small（#1）}}}

\begin{document}
\pagestyle{plain}
\begin{center}
{\LARGE\bfseries 上海高考高分翻译背诵 \quad 53}\\[0.25em]
{\large 每日背诵-53　2026-08-13 \quad\ 默写（中 $\to$ 英）}\\[0.35em]
{\footnotesize 来源：微信公众号「发现之旅 速来记单词」· 2026-08-13 · 赵文瀚}
\end{center}

\vspace{0.4em}
\noindent\begin{tabular}{@{}p{0.33\textwidth}@{}p{0.34\textwidth}@{}p{0.33\textwidth}@{}}
姓名：\underline{\hspace{2.0cm}} & 日期：\underline{\hspace{2.0cm}} & 得分：\underline{\hspace{1.4cm}}\,/\enspace 5 \\
\end{tabular}

\vspace{0.5em}
\noindent{\small 根据中文句子与提示词写出英文翻译，题号沿用原文。}

\vspace{0.8em}
\num{1}\textbf{人们倾向於无视这些抱怨。}

\vspace{0.3em}

\wline\vspace{0.3em}
\wline

\num{2}\textbf{音乐环绕着我，将凡俗事物隔离在外。}

\vspace{0.3em}

\wline\vspace{0.3em}
\wline

\num{3}\textbf{他向来衣着整洁素雅，合乎他的年龄和身份。}

\vspace{0.3em}

\wline\vspace{0.3em}
\wline

\num{4}\textbf{若漠视生活中的小确幸，你的世界将会黯然失色，枯燥乏味。}\prompt{or}

\vspace{0.3em}

\wline\vspace{0.3em}
\wline

\num{5}\textbf{近年来，该地区牛奶消耗量稳步提升，足见民众的健康意识已然深入人心。}\prompt{which}

\vspace{0.3em}

\wline\vspace{0.3em}
\wline

\end{document}
